% Options for packages loaded elsewhere
\PassOptionsToPackage{unicode}{hyperref}
\PassOptionsToPackage{hyphens}{url}
\documentclass[
]{article}
\usepackage{xcolor}
\usepackage{amsmath,amssymb}
\setcounter{secnumdepth}{-\maxdimen} % remove section numbering
\usepackage{iftex}
\ifPDFTeX
  \usepackage[T1]{fontenc}
  \usepackage[utf8]{inputenc}
  \usepackage{textcomp} % provide euro and other symbols
\else % if luatex or xetex
  \usepackage{unicode-math} % this also loads fontspec
  \defaultfontfeatures{Scale=MatchLowercase}
  \defaultfontfeatures[\rmfamily]{Ligatures=TeX,Scale=1}
\fi
\usepackage{lmodern}
\ifPDFTeX\else
  % xetex/luatex font selection
\fi
% Use upquote if available, for straight quotes in verbatim environments
\IfFileExists{upquote.sty}{\usepackage{upquote}}{}
\IfFileExists{microtype.sty}{% use microtype if available
  \usepackage[]{microtype}
  \UseMicrotypeSet[protrusion]{basicmath} % disable protrusion for tt fonts
}{}
\makeatletter
\@ifundefined{KOMAClassName}{% if non-KOMA class
  \IfFileExists{parskip.sty}{%
    \usepackage{parskip}
  }{% else
    \setlength{\parindent}{0pt}
    \setlength{\parskip}{6pt plus 2pt minus 1pt}}
}{% if KOMA class
  \KOMAoptions{parskip=half}}
\makeatother
\setlength{\emergencystretch}{3em} % prevent overfull lines
\providecommand{\tightlist}{%
  \setlength{\itemsep}{0pt}\setlength{\parskip}{0pt}}
\usepackage{bookmark}
\IfFileExists{xurl.sty}{\usepackage{xurl}}{} % add URL line breaks if available
\urlstyle{same}
\hypersetup{
  pdftitle={Cryptography Fundamentals},
  hidelinks,
  pdfcreator={LaTeX via pandoc}}

\title{Cryptography Fundamentals}
\author{}
\date{}

\begin{document}
\maketitle

\subsection{Description}\label{description}

Cryptography is the foundation of modern information security. This
course provides a comprehensive introduction to cryptographic concepts,
algorithms, and protocols. Topics include symmetric and asymmetric
encryption, hash functions, digital signatures, key exchange,
\texttt{TLS}/\texttt{SSL}, \texttt{PKI}, password storage, and
post-quantum cryptography. Participants will gain both theoretical
understanding and practical experience using libraries such as
\texttt{OpenSSL}, libsodium, and NaCl.

\subsection{Duration}\label{duration}

24 hours / 3 days

\subsection{Intended Audience}\label{intended-audience}

\begin{itemize}
\tightlist
\item
  Software developers implementing security features
\item
  Security engineers and analysts
\item
  Backend developers working with authentication and encryption
\item
  \texttt{DevOps} engineers managing certificates and secrets
\item
  Anyone who needs to understand cryptography for their technical role
\end{itemize}

\subsection{Prerequisites}\label{prerequisites}

\begin{itemize}
\tightlist
\item
  Proficiency in at least one programming language (\texttt{Python}, C,
  \texttt{Java}, or similar)
\item
  Basic understanding of computer science fundamentals (binary, data
  structures)
\item
  Familiarity with networking concepts (\texttt{TCP}/\texttt{IP},
  \texttt{HTTP}/\texttt{HTTPS})
\item
  Basic knowledge of security concepts (confidentiality, integrity,
  authentication)
\item
  Comfort with command-line tools
\end{itemize}

\subsection{Required Knowledge}\label{required-knowledge}

\begin{itemize}
\tightlist
\item
  \texttt{Python} Programming (or equivalent experience)
\end{itemize}

\subsection{Objectives}\label{objectives}

\begin{itemize}
\tightlist
\item
  Understand core cryptographic concepts and terminology
\item
  Apply symmetric encryption algorithms (\texttt{AES}, ChaCha20) with
  appropriate modes
\item
  Understand and use asymmetric encryption (\texttt{RSA}, ECC,
  Diffie-Hellman)
\item
  Use hash functions (\texttt{SHA}-256, \texttt{SHA}-3, BLAKE) and MACs
  for integrity verification
\item
  Implement and verify digital signatures
\item
  Understand the \texttt{TLS}/\texttt{SSL} protocol and certificate
  management
\item
  Apply proper password storage techniques (bcrypt, scrypt, Argon2)
\item
  Recognize and avoid common cryptographic mistakes
\item
  Use cryptographic libraries (\texttt{OpenSSL}, libsodium, NaCl)
  correctly
\item
  Understand the landscape of post-quantum cryptography
\end{itemize}

\subsection{Outline}\label{outline}

\begin{itemize}
\tightlist
\item
  Cryptography Concepts

  \begin{itemize}
  \tightlist
  \item
    History and evolution of cryptography
  \item
    Security goals: confidentiality, integrity, authentication,
    non-repudiation
  \item
    Kerckhoffs' principle and why secrecy of the algorithm is not
    security
  \item
    Computational hardness and security parameters
  \item
    Cryptographic primitives overview
  \item
    Standards bodies and cryptographic recommendations (\texttt{NIST},
    IETF)
  \end{itemize}
\item
  Symmetric Encryption

  \begin{itemize}
  \tightlist
  \item
    Symmetric encryption fundamentals and key management
  \item
    \texttt{AES} (Advanced Encryption Standard) algorithm and key sizes
  \item
    ChaCha20 stream cipher and its advantages
  \item
    \texttt{When} to use \texttt{AES} vs ChaCha20
  \item
    Key generation and secure key storage
  \end{itemize}
\item
  Block Cipher Modes

  \begin{itemize}
  \tightlist
  \item
    ECB (Electronic Codebook) and why it is insecure
  \item
    CBC (Cipher Block Chaining) mode and initialization vectors
  \item
    CTR (Counter) mode and parallelization
  \item
    GCM (Galois/Counter Mode) for authenticated encryption
  \item
    Choosing the right mode for your use case
  \item
    Padding schemes and padding \texttt{oracle} attacks
  \end{itemize}
\item
  Asymmetric Encryption

  \begin{itemize}
  \tightlist
  \item
    Public key cryptography fundamentals
  \item
    \texttt{RSA} algorithm, key generation, and key sizes
  \item
    Elliptic Curve Cryptography (ECC) and curve selection
  \item
    Diffie-Hellman key exchange protocol
  \item
    Elliptic Curve Diffie-Hellman (\texttt{ECDH})
  \item
    Hybrid encryption (combining symmetric and asymmetric)
  \item
    Performance comparison of asymmetric algorithms
  \end{itemize}
\item
  Hash Functions

  \begin{itemize}
  \tightlist
  \item
    Cryptographic hash function properties (collision resistance,
    preimage resistance)
  \item
    \texttt{SHA}-256 and the \texttt{SHA}-2 family
  \item
    \texttt{SHA}-3 (Keccak) and its design
  \item
    BLAKE2 and BLAKE3 for high-performance hashing
  \item
    Hash function use cases (integrity, fingerprinting, commitments)
  \item
    Birthday attacks and hash length considerations
  \end{itemize}
\item
  MACs and \texttt{HMAC}

  \begin{itemize}
  \tightlist
  \item
    Message Authentication Codes concept
  \item
    \texttt{HMAC} construction and security properties
  \item
    CMAC and GMAC alternatives
  \item
    Authenticated encryption with associated data (AEAD)
  \item
    \texttt{AES}-GCM and ChaCha20-Poly1305 as AEAD constructions
  \item
    \texttt{When} to use MACs vs digital signatures
  \end{itemize}
\item
  Digital Signatures

  \begin{itemize}
  \tightlist
  \item
    Digital signature concepts and use cases
  \item
    \texttt{RSA} signatures (\texttt{PKCS\#1\ v1.5}, PSS)
  \item
    \texttt{ECDSA} (Elliptic Curve Digital Signature Algorithm)
  \item
    EdDSA and Ed25519 signatures
  \item
    Signature verification and trust models
  \item
    Code signing and document signing applications
  \end{itemize}
\item
  Key Exchange Protocols

  \begin{itemize}
  \tightlist
  \item
    Key exchange problem and man-in-the-middle attacks
  \item
    Diffie-Hellman key exchange in detail
  \item
    Authenticated key exchange protocols
  \item
    Forward secrecy and ephemeral keys
  \item
    Key derivation functions (\texttt{HKDF}, \texttt{PBKDF2})
  \item
    Key agreement vs key transport
  \end{itemize}
\item
  \texttt{TLS}/\texttt{SSL} Protocol

  \begin{itemize}
  \tightlist
  \item
    \texttt{TLS} protocol overview and version history
  \item
    \texttt{TLS\ 1.3} handshake and improvements over \texttt{TLS\ 1.2}
  \item
    Cipher suite negotiation and selection
  \item
    Certificate validation during handshake
  \item
    Session resumption and 0-RTT
  \item
    Common \texttt{TLS} misconfigurations and vulnerabilities
  \item
    Testing \texttt{TLS} configuration
  \end{itemize}
\item
  \texttt{PKI} and Certificates

  \begin{itemize}
  \tightlist
  \item
    Public Key Infrastructure concepts
  \item
    \texttt{X.509} certificate structure and fields
  \item
    Certificate Authorities (\texttt{CA}) and trust chains
  \item
    Certificate signing requests (CSR)
  \item
    Certificate revocation (CRL, \texttt{OCSP})
  \item
    \texttt{Let\textquotesingle{}s\ Encrypt} and automated certificate
    management
  \item
    Certificate pinning and its trade-offs
  \end{itemize}
\item
  Key Management

  \begin{itemize}
  \tightlist
  \item
    Key lifecycle (generation, distribution, storage, rotation,
    destruction)
  \item
    Hardware Security Modules (\texttt{HSM})
  \item
    Key Management Services (\texttt{cloud\ KMS})
  \item
    Key wrapping and envelope encryption
  \item
    Secret sharing schemes
  \item
    Key escrow and recovery
  \end{itemize}
\item
  Password Storage

  \begin{itemize}
  \tightlist
  \item
    Why simple hashing is insufficient
  \item
    Salting and its importance
  \item
    bcrypt algorithm and cost factor tuning
  \item
    scrypt and memory-hard functions
  \item
    Argon2 (Argon2id) and \texttt{OWASP} recommendations
  \item
    Password hashing benchmarks and parameter selection
  \item
    Credential stuffing defense
  \end{itemize}
\item
  Random Number Generation

  \begin{itemize}
  \tightlist
  \item
    Importance of randomness in cryptography
  \item
    Pseudorandom number generators (PRNG) vs cryptographically secure
    PRNGs (CSPRNG)
  \item
    Operating system entropy sources (/dev/urandom, CryptGenRandom)
  \item
    Common mistakes with random number generation
  \item
    Testing randomness quality
  \end{itemize}
\item
  Post-Quantum Cryptography

  \begin{itemize}
  \tightlist
  \item
    Quantum computing threat to current cryptography
  \item
    \texttt{NIST} post-quantum standardization process
  \item
    Lattice-based cryptography (CRYSTALS-Kyber, CRYSTALS-Dilithium)
  \item
    Hash-based signatures (SPHINCS+)
  \item
    Migration planning and crypto agility
  \item
    Hybrid classical/post-quantum approaches
  \end{itemize}
\item
  Common Cryptographic Mistakes

  \begin{itemize}
  \tightlist
  \item
    Rolling your own crypto
  \item
    Using ECB mode for encryption
  \item
    Reusing nonces and initialization vectors
  \item
    Insufficient key lengths and weak algorithms
  \item
    Improper random number generation
  \item
    Timing attacks and side-channel vulnerabilities
  \item
    Not using authenticated encryption
  \item
    Hardcoded keys and secrets
  \end{itemize}
\item
  Cryptographic Libraries

  \begin{itemize}
  \tightlist
  \item
    \texttt{OpenSSL} command-line and library usage
  \item
    libsodium high-level \texttt{API} and best practices
  \item
    NaCl (Networking and Cryptography Library) design philosophy
  \item
    Language-specific libraries (\texttt{Python} cryptography,
    \texttt{Java} JCA/JCE, Go crypto)
  \item
    Choosing the right library for your project
  \item
    Keeping libraries updated and monitoring advisories
  \end{itemize}
\end{itemize}

\subsection{Copyright}\label{copyright}

Mark Veltzer
\href{mailto:mark.veltzer@gmail.com}{\nolinkurl{mark.veltzer@gmail.com}},
© 2026

\end{document}
